\section{Scope: two independent supersymmetry bricks}

The aim is deliberately modular. We isolate two exact supersymmetry constructions and keep their assumptions separate.

\begin{enumerate}[label=\textbf{S\arabic*.}]
\item An operator-index brick: Fredholm factorization, Toeplitz winding, Hilbert-Hotel shifts, and protected zero modes in $\mathcal N=2$ supersymmetric quantum mechanics.
\item A tetrahedral-spinor brick: a qubit SIC/Pauli orbit, its binary-tetrahedral spin symmetry, and an $\mathcal N=1$ supertranslation realization after explicit CAR input.
\end{enumerate}

No theorem requires the two bricks to be identified. Their value is that each can be checked, extended, or falsified independently.

\subsection{Euler/Pauli grading}

Euler's half-turn
\[
e^{i\pi}=-1
\]
appears in the spinorial lift as a central sign. For any Pauli axis,
\[
U_j(\pi)=\exp\!\left(-\frac{i\pi}{2}\sigma_j\right)=-i\sigma_j,
\qquad
U_j(\pi)^2=-I,
\qquad
U_j(\pi)^4=I.
\]
Thus an order-two projective half-turn lifts to an order-four spin operation.

The companion result SOH-G013 proves the determinant-one Pauli lifts
\[
\widetilde R=i\sigma_x,
\qquad
\widetilde E=i\sigma_z,
\qquad
\widetilde N=i\sigma_y,
\]
which generate the quaternion group $Q_8$ and project to the Klein four group $V_4$. The grading used below is
\[
\Gamma=\sigma_z\otimes I,
\qquad
\Gamma^2=I.
\]
The key strengthening in the present revision is that Euler phase also enters through winding: the same $U(1)$ phase variable that supplies the half-turn carries an integer topological charge on the boundary circle.
